%!TEX TS-program = xelatex
\documentclass[11pt, a4paper]{article}

\usepackage[fontset=windows]{ctex}
\usepackage{amsmath, amssymb, amsthm}
\usepackage{graphicx}
\usepackage{hyperref}
\usepackage{bookmark}
\usepackage{enumitem}

\newtheorem{problem}{题}
\renewcommand{\proofname}{\textbf{解答}}
\usepackage[margin=1in]{geometry}

\title{算法设计与分析 作业七}
\author{xxxxxxxxx xxx}
\date{\today}

\begin{document}

\maketitle

\begin{problem}[9.5]
\end{problem}

\begin{proof}
    \begin{enumerate}[label=(\arabic*)]
        \item 由独立集到团。给定无向图 $G=(V,E)$ 和整数 $J$，构造补图
        $\overline{G}=(V,\overline{E})$，其中
        \[
            \overline{E}=\{(u,v)\mid u,v\in V,\ u\ne v,\ (u,v)\notin E\}.
        \]
        并令团问题中的整数仍为 $J$。若 $S$ 是 $G$ 中大小至少为 $J$ 的独立集，则任意
        $u,v\in S$ 在 $G$ 中不相邻，于是在 $\overline{G}$ 中相邻，所以 $S$ 是
        $\overline{G}$ 中的团。反过来也成立。因此该变换是多项式时间变换。

        \item 由顶点覆盖到团。给定无向图 $G=(V,E)$ 和整数 $K$，构造补图
        $\overline{G}$，并令
        \[
            J=|V|-K.
        \]
        若 $C$ 是 $G$ 中大小不超过 $K$ 的顶点覆盖，则 $V-C$ 是 $G$ 的独立集，因而是
        $\overline{G}$ 中大小至少为 $J$ 的团。反过来，若 $\overline{G}$ 有大小至少为
        $J$ 的团 $S$，则 $S$ 是 $G$ 的独立集，所以 $V-S$ 是 $G$ 的顶点覆盖，且大小不超过
        $K$。故 $VC\le_p$ 团。

        \item 由哈密顿回路到子图同构。给定无向图 $G=(V,E)$，令
        $H=(V_1,E_1)$ 为含 $|V|$ 个顶点的简单环。若 $H$ 同构于 $G$ 的某个子图，则这个
        子图正是一条经过所有顶点一次的回路，即 $G$ 有哈密顿回路；反之，$G$ 的哈密顿
        回路显然给出一个与 $H$ 同构的子图。因此 $HC\le_p$ 子图同构。

        \item 由团到 $0$-$1$ 整数规划。给定无向图 $G=(V,E)$ 和整数 $J$，设
        $V=\{v_1,\ldots,v_n\}$，为每个顶点引入变量 $x_i\in\{0,1\}$。构造约束
        \[
            \sum_{i=1}^n x_i\ge J,
        \]
        并且对每一对不相邻的顶点 $v_i,v_j$ 加入
        \[
            x_i+x_j\le 1.
        \]
        若 $G$ 中存在大小至少为 $J$ 的团，则把团中顶点对应变量取 $1$，其余取 $0$，得到
        可行解。反之，任何可行解中所有取 $1$ 的顶点两两相邻，且个数至少为 $J$，所以它们
        构成团。该整数规划规模显然是多项式的。
    \end{enumerate}
\end{proof}

\vspace{1em}

\begin{problem}[9.10]
\end{problem}

\begin{proof}
    已知 $\Pi'$ 是 NP 完全问题，且 $\Pi'\le_p \Pi$。由于 $\Pi'$ 是 NP 完全的，对任意
    $\Pi''\in NP$，都有
    \[
        \Pi''\le_p \Pi'.
    \]
    又因为多项式时间归约具有传递性，得到
    \[
        \Pi''\le_p \Pi'\le_p \Pi,
    \]
    从而任意 NP 问题都可多项式时间归约到 $\Pi$，即 $\Pi$ 是 NP 难的。

    题设还给出 $\Pi\in NP$，所以 $\Pi$ 既属于 NP 又是 NP 难的，故 $\Pi$ 是 NP 完全的。
\end{proof}

\vspace{1em}

\begin{problem}[9.15]
\end{problem}

\begin{proof}
    先说明子图同构属于 NP。证书可以取为一个映射
    $f:V_1\to V$。验证时只需检查 $f$ 是否为单射，并对每条
    $(u,v)\in E_1$ 检查 $(f(u),f(v))\in E$，这些都能在多项式时间内完成。

    再证明它是 NP 难的。由哈密顿回路问题归约。给定哈密顿回路实例
    $G=(V,E)$，构造子图同构实例 $(G,H)$，其中 $H$ 是一个包含 $|V|$ 个顶点的简单环。
    若 $G$ 有哈密顿回路，则这条回路给出了 $H$ 到 $G$ 的子图同构。反过来，若 $G$ 中
    有一个子图与 $H$ 同构，因为 $H$ 的顶点数正好为 $|V|$，该子图必经过 $G$ 的所有顶点，
    因而就是 $G$ 的哈密顿回路。

    该构造显然可在多项式时间内完成，且哈密顿回路是 NP 完全问题，所以子图同构是
    NP 难的。结合子图同构属于 NP，故子图同构是 NP 完全的。
\end{proof}

\vspace{1em}

\begin{problem}[9.19]
\end{problem}

\begin{proof}
    先说明支配集属于 NP。证书为一个顶点子集 $V'\subseteq V$。验证时检查
    $|V'|\le K$，并检查每个顶点要么属于 $V'$，要么与 $V'$ 中某个顶点相邻，均可在
    多项式时间内完成。

    下面由顶点覆盖归约到支配集。给定顶点覆盖实例 $G=(V,E)$ 和整数 $K$，构造图
    $G'=(V',E')$。令
    \[
        V'=V\cup E,
    \]
    即把原图每条边也看成一个新顶点。边集由两部分组成：第一部分使原顶点集合 $V$ 构成
    完全图；第二部分中，对每条原边 $e=(u,v)$，把新顶点 $e$ 分别与 $u$ 和 $v$ 相连。令
    支配集实例的整数仍为 $K$。该构造规模为多项式。

    若 $C$ 是 $G$ 中大小不超过 $K$ 的顶点覆盖，则 $C$ 在 $G'$ 中支配所有顶点：原顶点
    由于 $V$ 构成完全图而被 $C$ 支配；每个边顶点 $e=(u,v)$ 至少有一个端点在 $C$ 中，
    因而也被支配。所以 $G'$ 有大小不超过 $K$ 的支配集。

    反过来，设 $D$ 是 $G'$ 中大小不超过 $K$ 的支配集。若 $D$ 中含有边顶点
    $e=(u,v)$，则可用端点 $u$ 替换它。因为 $V$ 是完全图，$u$ 能支配所有原顶点；同时
    $u$ 至少支配边顶点 $e$，替换不会增加集合大小。不断替换后，可得到一个只含原顶点的
    支配集 $D'\subseteq V$，且 $|D'|\le K$。若原图某条边 $e=(u,v)$ 的两个端点都不在
    $D'$ 中，则边顶点 $e$ 在 $G'$ 中只与 $u,v$ 相邻，且自身也不在 $D'$ 中，无法被支配，
    矛盾。因此 $D'$ 覆盖了 $G$ 的每条边，是大小不超过 $K$ 的顶点覆盖。

    于是顶点覆盖多项式时间归约到支配集。由于顶点覆盖是 NP 完全的，支配集属于 NP 且
    NP 难，所以支配集是 NP 完全的。
\end{proof}

\end{document}
